\section{Polynomial lifting through the actual residue-duality
comparison}\label{polynomial-lifting-through-the-actual-residue-duality-comparison}

24 September 2026. Independent derivation DPL0--DPL9.

\subsection{DPL0. Actual source, prior results, and the receiving
category}\label{dpl0.-actual-source-prior-results-and-the-receiving-category}

This note uses the original quotient and the actual comparison cone
constructed in
\href{GLOBAL_ORIGINAL_ZETA_RESIDUE_DUALITY_INDEPENDENT.md}{GZR} and
\href{ACTUAL_SUPPORTED_DUALITY_INDEPENDENT.md}{ASD14}. It proves exact
polynomial lifting and extension consequences for their entire dual
cokernel. It does not assume that this cokernel is zero.

The source notation remains \(Z_0,Z_1/\tau,Z_2\). The current operations
and retractions in \nolinkurl{CORPUS_AND_OPERATION_RULES.md} remain controlling.
The complete global user arguments USR-9ad1c0a2d09dba92,
USR-f55d16feb948d8c2, and USR-6152e3bc6302258c were read in the
preceding GZR derivation and are not replaced by a new source definition
here. Complex addition, polynomial operators, and exact sequences below
belong to the already constructed receiving spaces. Primitive \(\tau\)
is not an operand of any added operation.

The exact prior inputs are
\href{ORIGINAL_MELLIN_SPECTRAL_SYNTHESIS.md}{OMS1--OMS6}, the original
test-space and quotient identification;
\href{https://github.com/KokunoYumeto/zeta-function-research-reader/blob/main/workbenches/splitzero-tandem/continuations/20260924-cohomology-weight-and-character-lifts/gct-weight-control/proofs/ORIGINAL_ZETA_RESOLVENT_AND_PROJECTORS.md}{RZ2--RZ8},
the continuous resolvent and actual full-jet representatives;
GZR3--GZR8, the original-zeta residue pairing; and ASD2, ASD4, and
ASD14, the continuous-dual topology and the actual comparison cone.
ASD14 was read completely for this calculation. The original analytic
source remains Ralf Meyer,
\href{https://arxiv.org/abs/math/0412277v3}{arXiv:math/0412277v3}, as
entered with full original factors in OMS. The receiving coefficient
complex is that of Alain Connes and Caterina Consani,
\href{https://arxiv.org/abs/0903.2024v3}{arXiv:0903.2024v3, §5}, through
the exact ASD source comparison. No new reading of these entire human
papers is claimed here.

Retain the entire-function Fréchet algebra \[
\mathcal B=\{F\text{ entire}:\
b_{A,N}(F)=\sup_{|\sigma|\le A,t\in\mathbb R}
(1+|t|)^N|F(\sigma+it)|<\infty\text{ for all }A,N\},
\] \[
\mathcal I=\{F\in\mathcal B:F^{(j)}(\rho)=0\
(\rho\in\mathscr Z,\ 0\le j<m_\rho)\},\qquad Q=\mathcal B/\mathcal I.
\tag{DPL0.1}
\] Here \(\mathscr Z\) is the actual nontrivial zero set of the original
\(\zeta\), with every actual multiplicity. OMS proves that this quotient
is the original \(\mathcal A/\mathcal ES\), with its quotient topology.
Write \(L[F]=[sF]\), \(T_a[F]=[a^sF]\) for \(a>0\), and \(Q'\) for the
continuous complex-linear dual.

The constructed pairing and its injective dual map are \[
B_\zeta([F],[G])=
\frac1{2\pi i}\left(\int_{2-i\infty}^{2+i\infty}
-\int_{-1-i\infty}^{-1+i\infty}\right)
\frac{F(s)G(1-s)}{\zeta(s)}\,ds,\qquad
D x(y)=B_\zeta(x,y).
\tag{DPL0.2}
\] GZR proves both integrals converge absolutely, the form descends in
both slots, every full primary residue is retained, and \(D:Q\to Q'_b\)
is continuous and injective. Here \(Q'_b\) denotes the strong continuous
dual. ASD also uses the weak-* dual \(Q'_\sigma\).

Write \[
A=L^t,\qquad (A\lambda)(x)=\lambda(Lx),\qquad
C=Q'/DQ,\qquad \eta:Q'\to C.
\tag{DPL0.3}
\] The algebraic quotient is retained, whether or not its quotient
topology is Hausdorff. We use either the specified weak-* quotient
topology or, when explicitly named, the quotient of the strong dual.
These topologies are not identified with each other.

The exact real-action identity and its infinitesimal form are \[
DT_a=aT_a^\vee D,\qquad
DL=(1-A)D,\qquad
(T_a^\vee\lambda)(x)=\lambda(T_{a^{-1}}x).
\tag{DPL0.4}
\] Thus the actual twisted target \(Y=\chi\otimes Q'\), \(\chi(a)=a\),
has generator \[
M=1-A,
\tag{DPL0.5}
\] and \(D:(Q,L)\to(Y,M)\) is an intertwiner. The twisted cokernel
\(\widetilde C=\chi\otimes C\) has the same vectors as \(C\), with
induced generator \(1-A\). This twist is kept throughout.

For Hom and Ext statements the receiving module ring is
\(R=\mathbb C[X]\), with \(X\) acting by the specified generator. Thus
\(X\) acts by \(L\) on \(Q\) and by \(M\) on \(Y,\widetilde C\). These
are the actual modules from (DPL0.1)--(DPL0.5), not newly postulated
spectra. Algebraic \(\operatorname{Ext}^1_R\) means equivalence classes
of exact module extensions. A continuous assertion is stated for strict
exact rows of locally convex receiving modules: the injection has its
subspace topology, and the surjection has its quotient topology.
Non-Hausdorff quotients are not silently replaced by their
Hausdorffizations.

The phrase ``receiving polynomial torsion'' below means that each vector
in this \(R\)-module is annihilated by some nonzero polynomial in its
specified receiving generator. It does not mean the user's
integer-amount datum, which the corpus sometimes also calls torsion. A
single polynomial annihilating the entire module is a stronger,
separately stated condition.

\subsection{DPL1. Explicit continuous division on the original
quotient}\label{dpl1.-explicit-continuous-division-on-the-original-quotient}

Let \(a\in\mathbb C\), and write \(S_a=L-a\). If \(a\notin\mathscr Z\),
the original RZ resolvent proves that \(S_a\) is a continuous
automorphism of \(Q\). Its actual representative inverse is \[
\mathscr V_aF(s)=
\frac{F(s)-F_*(s)F(a)/F_*(a)}{s-a},
\qquad
F_*(s)=\frac{s(s-1)}8\pi^{-s/2}\Gamma(s/2)\zeta(s).
\tag{DPL1.1}
\] All endpoint and trivial-zero values in RZ1 remain part of this
formula. In particular \(F_*(0)=F_*(1)=1/8\); division at either
endpoint is not division by a removed original zeta value.

Now let \(a\in\mathscr Z\), \(m=m_a\). Let \(e_a=E_{a,0}\) be the actual
RZ representative with Taylor jet one at \(a\) and zero to the full
required order at every other actual zero. Set \[
E_{a,j}(s)=e_a(s)(s-a)^j,\quad 0\le j<m,\qquad
\ell_{a,j}([F])=F^{(j)}(a)/j!.
\tag{DPL1.2}
\] All \(\ell_{a,j}\) with \(j<m\) are well-defined and continuous on
\(Q\). Define on the entire original test algebra \[
\mathscr U_aF(s)=
\frac{F(s)-F(a)e_a(s)}{s-a}
-\left(\frac{F^{(m)}(a)}{m!}
-F(a)\frac{e_a^{(m)}(a)}{m!}\right)E_{a,m-1}(s).
\tag{DPL1.3}
\] The first numerator vanishes at \(a\). Its entire division is
continuous in \(\mathcal B\): outside a fixed disk around \(a\), divide
by the bounded-away-from-zero linear factor; inside the disk use the
derivative integral along the segment from \(a\), with Cauchy's estimate
on a larger disk. This bounds each output strip seminorm by finitely
many input strip seminorms. Every remaining operation in (DPL1.3) is a
continuous finite-rank operation.

If \(F\in\mathcal I\), then \(F(a)=0\), \(F^{(j)}(a)=0\) for \(j<m\),
and the first quotient has exactly the possible last jet
\(F^{(m)}(a)t^{m-1}/m!\). The explicit correction removes that last jet.
At every other zero both terms retain its full vanishing order. Thus
\(\mathscr U_a\mathcal I\subset\mathcal I\), and (DPL1.3) induces a
continuous \(U_a:Q\to Q\).

Its identities are \[
S_aU_a=1-[E_{a,0}]\ell_{a,0},\qquad
U_aS_a=1-[E_{a,m-1}]\ell_{a,m-1}.
\tag{DPL1.4}
\] For the first, multiplying (DPL1.3) by \(s-a\) leaves \(F-F(a)e_a\)
modulo \(\mathcal I\), since \((s-a)E_{a,m-1}\in\mathcal I\). For the
second, replace \(F\) in (DPL1.3) by \((s-a)F\): its value at \(a\) is
zero and its coefficient of degree \(m\) is \(F^{(m-1)}(a)/(m-1)!\).
This proves the second identity even at the representative level.

Consequently \[
\ker S_a=\mathbb C[E_{a,m-1}],\qquad
\operatorname{im}S_a=\ker\ell_{a,0}.
\tag{DPL1.5}
\] The first identity follows by applying \(U_aS_a\) to a kernel vector
and observing that \(S_a[E_{a,m-1}]=0\). For the second, evaluation at
\(a\) kills the range, and \(S_aU_a\) is identity on \(\ker\ell_{a,0}\).
The range is closed. The displayed continuous right inverse on the range
and continuous projection away from the kernel give strict division on
the actual quotient, with every multiplicity retained.

Transposing (DPL1.4) gives \[
(A-a)U_a^t
=1-\ell_{a,m-1}\operatorname{ev}_{[E_{a,m-1}]},
\] \[
U_a^t(A-a)
=1-\ell_{a,0}\operatorname{ev}_{[E_{a,0}]}.
\tag{DPL1.6}
\] Here \(\operatorname{ev}_v(\lambda)=\lambda(v)\). Therefore \[
\ker(A-a)=\mathbb C\ell_{a,0},\qquad
\operatorname{im}(A-a)=\{\lambda:\lambda([E_{a,m-1}])=0\}.
\tag{DPL1.7}
\] These identities follow directly from (DPL1.6); no extension of a
discontinuous functional is involved. Each transpose is continuous for
both dual topologies because continuous operators preserve bounded sets
and weak-* evaluations.

\subsection{DPL2. Every linear factor is invertible on the actual dual
cokernel}\label{dpl2.-every-linear-factor-is-invertible-on-the-actual-dual-cokernel}

First \(DQ\) is \(A\)-stable by (DPL0.4), so \(A-a\) induces an operator
on \(C\). For \(a\notin\mathscr Z\), also \(1-a\notin\mathscr Z\) by the
original functional equation. Both source operators in \[
(A-a)D=-D(L-(1-a))
\tag{DPL2.1}
\] are continuously invertible on the appropriate full spaces. The
transpose inverse on \(Q'\) therefore induces an inverse on \(C\).

For \(a\in\mathscr Z\), put \(b=1-a\), \(m=m_a=m_b\), and retain \[
\zeta(b+t)=t^mu_b(t),\qquad
c_{b,0}=u_b(0)^{-1}=\frac{m!}{\zeta^{(m)}(b)}.
\tag{DPL2.2}
\] The finite-residue theorem GZR5, on the original denominator, gives
the two exact identities \[
D[E_{b,m-1}]=c_{b,0}\ell_{a,0},
\qquad
D[E_{b,0}]([E_{a,m-1}])=(-1)^{m-1}c_{b,0}.
\tag{DPL2.3}
\] For the first identity evaluate on an arbitrary \(F\): only the
residue at \(b\) remains, and its coefficient is \(c_{b,0}F(a)\). For
the second, the reflected monomial contributes \((-1)^{m-1}\), with no
other allowed residue term. These calculations retain the actual
derivative \(\zeta^{(m)}(b)\).

To prove surjectivity on \(C\), take any \(\lambda\in Q'\), and put \[
h_\lambda=
\frac{\lambda([E_{a,m-1}])}{(-1)^{m-1}c_{b,0}}[E_{b,0}],
\qquad
K_a\lambda=U_a^t(\lambda-Dh_\lambda).
\tag{DPL2.4}
\] The adjusted functional annihilates \([E_{a,m-1}]\). Equation
(DPL1.6) therefore proves \[
(A-a)K_a\lambda=\lambda-Dh_\lambda.
\tag{DPL2.5}
\] Taking classes gives every prescribed class as an image.

To prove injectivity, suppose \((A-a)\lambda=Dh\). Evaluation on
\([E_{a,m-1}]\) gives zero on the left. By the same full residue
calculation, the right side is \((-1)^{m-1}c_{b,0}\ell_{b,0}(h)\). Thus
\(\ell_{b,0}(h)=0\). Put \(k=-U_bh\); then \((L-b)k=-h\) by (DPL1.4).
Equation (DPL2.1) gives \[
(A-a)Dk=Dh=(A-a)\lambda.
\] Hence \(\lambda-Dk\in\ker(A-a)=\mathbb C\ell_{a,0}\). This kernel
lies in \(DQ\) by (DPL2.3), proving \(\lambda\in DQ\). The induced
kernel on \(C\) is zero.

Thus every \(A-a\) is bijective on \(C\). The explicit operator \(K_a\)
is continuous on \(Q'\) for either topology: it is a transpose composed
with identity minus a finite-rank evaluation map. Equation (DPL2.5) and
injectivity on \(C\) imply \(K_a(DQ)\subset DQ\), so it induces the
unique inverse on \(C\), continuously for each specified quotient
topology. The case \(a\notin\mathscr Z\) uses the transpose of the
continuous resolvent instead. Consequently every linear factor is a
continuous automorphism of the actual algebraic cokernel with either of
these quotient topologies.

\subsection{DPL3. Every polynomial and the retained character
twist}\label{dpl3.-every-polynomial-and-the-retained-character-twist}

A nonzero polynomial factors over \(\mathbb C\) as \[
q(X)=c\prod_{j=1}^r(X-a_j),\qquad c\ne0,
\tag{DPL3.1}
\] with repetitions retained. The commuting operators \(A-a_j\) are all
bijective on \(C\); their inverses commute because they invert commuting
bijections. Thus \[
q(A)^{-1}=c^{-1}\prod_{j=1}^r(A-a_j)^{-1}
\tag{DPL3.2}
\] is a continuous two-sided inverse. The empty product covers nonzero
constants.

On the actual twisted cokernel the generator is \(M=1-A\), so \[
q(M)=q(1-A)
\tag{DPL3.3}
\] is also continuously invertible. There is a unique action of the
rational-function field \(\mathbb C(X)\) on the actual cokernel given by
\[
\frac{p(X)}{q(X)}\cdot c=p(M)q(M)^{-1}c .
\tag{DPL3.4}
\] If \(p/q=p_1/q_1\), cross-multiplication and invertibility of
\(q(M)q_1(M)\) prove the two actions agree. The same identities prove
addition and multiplication of rational functions. This constructs an
action on the existing quotient, without asserting its dimension over
that field or that the quotient is nonzero.

In particular \[
\ker q(M)|_{\widetilde C}=0,\qquad
q(M)\widetilde C=\widetilde C\qquad(q\ne0).
\tag{DPL3.5}
\] This excludes every nonzero polynomial-annihilated vector in the
actual cokernel. It does not identify the entire cokernel with its
Hausdorff quotient or assert surjectivity of \(D\).

\subsection{DPL4. Both Hom vanishings and an explicit unique extension
splitting}\label{dpl4.-both-hom-vanishings-and-an-explicit-unique-extension-splitting}

Let \(V\) be an \(R\)-module with \(q(T_V)V=0\) for one specified
nonzero polynomial \(q\). Its vector space may be infinite-dimensional;
the hypothesis means this one polynomial annihilates the whole module.

For an \(R\)-linear \(f:V\to\widetilde C\), \[
q(M)f(v)=f(q(T_V)v)=0,
\] so injectivity of \(q(M)\) gives \(f=0\). For an \(R\)-linear
\(g:\widetilde C\to V\), every \(c\) has the form \(q(M)c_0\), giving
\(g(c)=q(T_V)g(c_0)=0\). Therefore \[
\boxed{\operatorname{Hom}_R(V,\widetilde C)=0
=\operatorname{Hom}_R(\widetilde C,V).}
\tag{DPL4.1}
\]

Consider any exact \(R\)-module extension \[
0\longrightarrow\widetilde C\xrightarrow{i}E
\xrightarrow{p}V\longrightarrow0.
\tag{DPL4.2}
\] Since \(p\,q(T_E)=q(T_V)p=0\), the operator \(q(T_E)\) has image in
\(i\widetilde C\). Define \[
r:E\to\widetilde C,\qquad
r(e)=q(M)^{-1}\,i^{-1}\bigl(q(T_E)e\bigr).
\tag{DPL4.3}
\] The inverse \(i^{-1}\) is taken only on its actual image. This is an
\(R\)-linear map: every operator in the formula commutes with the
indicated generator. Since \(q(T_E)i=i\,q(M)\), it satisfies \(ri=1\).

For \(v\in V\), choose any \(e\) with \(p(e)=v\), and set \[
\boxed{\sigma(v)=e-i\,r(e).}
\tag{DPL4.4}
\] Replacing \(e\) by \(e+i(c)\) changes this expression by
\(i(c)-ir i(c)=0\). Thus the formula is independent of the chosen lift.
It is linear and \(R\)-linear because \(1-ir\) is, and \(p\sigma=1\). It
is the unique \(R\)-linear splitting: the difference of two splittings
factors through \(i\) as an \(R\)-linear map \(V\to\widetilde C\), which
is zero by (DPL4.1). This proves \[
\boxed{\operatorname{Ext}^1_R(V,\widetilde C)=0.}
\tag{DPL4.5}
\]

For a strict exact row in the locally convex receiving category defined
in DPL0, (DPL4.3) is continuous: \(i^{-1}\) is continuous on its
subspace image, \(q(T_E)\) is continuous into that image, and DPL3
supplies the continuous inverse \(q(M)^{-1}\). Because \(p\) has its
quotient topology and \(\sigma p=1-ir\) is continuous, \(\sigma\) is
continuous as well. Thus the same explicit unique splitting proves the
corresponding strict continuous extension statement.

Any additional commuting receiving action is preserved by this
splitting. In particular, the original positive-real actions and the
source receiving ring actions commute with the generator; their
operators commute with \(q(M)\) and its inverse, so (DPL4.3)--(DPL4.4)
intertwine them. This is a full equivariance statement for extensions
that already carry the actual compatible actions, not an inferred action
on primitive \(\tau\).

The words ``one specified nonzero polynomial'' matter for the reverse
Hom statement. We have not replaced this condition by elementwise
torsion with no common polynomial. Forward Hom into \(\widetilde C\)
does vanish for elementwise torsion, because the same argument applies
separately to each vector; this stronger observation is not used to
enlarge the asserted reverse Hom theorem.

\subsubsection{DPL4A. The algebraic splitting for all receiving
polynomial-torsion
modules}\label{dpl4a.-the-algebraic-splitting-for-all-receiving-polynomial-torsion-modules}

There is a further exact algebraic extension result when the
annihilating polynomial depends on the vector. Let \(V\) have receiving
polynomial torsion in the sense just defined, and keep the algebraic
extension (DPL4.2). For each \(e\in E\), choose a nonzero polynomial
\(q\) annihilating \(p(e)\), and define \[
r_q(e)=q(M)^{-1}i^{-1}(q(T_E)e).
\tag{DPL4A.1}
\] This expression is independent of \(q\). Indeed if \(r\) is another
annihilating polynomial, multiplication of each proposed value by the
bijection \(q(M)r(M)\) gives the same vector \[
i^{-1}\bigl(q(T_E)r(T_E)e\bigr).
\] The equality follows because the two polynomials commute with each
other and with the injection. Injectivity of the product proves
independence.

For \(e_1,e_2\), the product of annihilating polynomials for their two
images also annihilates their sum. Using this single product in
(DPL4A.1) proves additivity of \(r\). The same polynomial handles scalar
multiplication, and a polynomial annihilating \(p(e)\) also annihilates
\(T_Vp(e)\); commutation proves \(r(T_Ee)=Mr(e)\). For \(e=i(c)\) one
may choose \(q=1\), giving \(r(i(c))=c\). Thus the independent formula
defines an \(R\)-linear retraction.

Formula (DPL4.4) now again gives a splitting. It is unique because an
\(R\)-linear map from receiving polynomial torsion into \(\widetilde C\)
vanishes vector by vector by (DPL3.5). Therefore \[
\operatorname{Hom}_R(V,\widetilde C)=0,\qquad
\operatorname{Ext}^1_R(V,\widetilde C)=0
\quad\text{for every receiving polynomial-torsion }V.
\tag{DPL4A.2}
\] This argument is algebraic. It does not supply a uniform continuity
estimate when the annihilating polynomial varies with the vector. The
continuous extension theorem in DPL4 retains its single global
polynomial and strict-topology hypotheses. Nor does (DPL4A.2) assert the
reverse Hom vanishing for arbitrary elementwise torsion.

\subsection{\texorpdfstring{DPL5. Every polynomial-annihilated dual
vector lifts uniquely to the actual
\(Q\)}{DPL5. Every polynomial-annihilated dual vector lifts uniquely to the actual Q}}\label{dpl5.-every-polynomial-annihilated-dual-vector-lifts-uniquely-to-the-actual-q}

For the actual exact algebraic row \[
0\longrightarrow Q\xrightarrow{D}Y=\chi\otimes Q'
\xrightarrow{\eta}\widetilde C\longrightarrow0,
\tag{DPL5.1}
\] take \(y\in Y\) with \(q(M)y=0\). Its image in \(\widetilde C\) is
killed by the same polynomial and hence is zero by DPL3. Thus \(y=Dx\)
for a unique \(x\in Q\). Intertwining gives \(Dq(L)x=q(M)y=0\), and
injectivity of \(D\) proves \(q(L)x=0\). Therefore \[
\boxed{
D:\ker q(L)|_Q\xrightarrow{\ \cong\ }\ker q(M)|_Y
}
\tag{DPL5.2}
\] for every nonzero polynomial, including every \((X-\rho)^r\).

Both spaces in (DPL5.2) are finite-dimensional with their actual
inherited topologies. To verify this on the source, factor \(q\). At an
actual zero \(\rho\), the polynomial germ has order
\(r_\rho=\operatorname{ord}_\rho q\), where \(r_\rho=0\) if
\(q(\rho)\ne0\). In the actual primary algebra
\(\mathbb C[t]/(t^{m_\rho})\), the kernel is precisely \[
t^{\max(m_\rho-r_\rho,0)}\mathbb C[t]/(t^{m_\rho}),
\tag{DPL5.3}
\] of dimension \(\min(m_\rho,r_\rho)\). Only finitely many zeros can
occur among the finitely many roots of \(q\). Every other primary jet is
zero. Summing the existing RZ primary sections reconstructs these
finitely many possible jets, and synthesis makes the remaining class
zero. Thus \[
\dim\ker q(L)=
\sum_{\rho:q(\rho)=0}\min(m_\rho,\operatorname{ord}_\rho q).
\tag{DPL5.4}
\] The same dimension holds on the target by (DPL5.2). These are
subspaces of Hausdorff locally convex spaces, so their
finite-dimensional topologies are the usual ones. The bijection in
(DPL5.2) and its inverse are continuous.

More generally, any \(R\)-linear map \(f:V\to Y\) from a module
annihilated by \(q\) has \(\eta f=0\) by DPL4.1, hence has a unique
\(R\)-linear lift \(\widetilde f:V\to Q\) satisfying
\(D\widetilde f=f\). Its image lies in the finite-dimensional kernel
(DPL5.3). If \(f\) is continuous, the finite-dimensional inverse just
proved makes \(\widetilde f\) continuous. If \(f\) also intertwines any
of the actual real or source-ring actions, so does its lift, by
injectivity and the exact intertwining of \(D\). No continuous inverse
of \(D\) on its entire image has been assumed.

There is an exact local formula for this lift, retaining every original
zeta derivative. Suppose \((A-a)^r\lambda=0\), let \(b=1-a\), and assume
\(a\in\mathscr Z\) of multiplicity \(m\). For its unique preimage \(x\)
under \(D\), the \(b\)-jet is \[
j_bx(t)=
\left[
u_b(t)\sum_{j=0}^{m-1}(-1)^j
\lambda([E_{a,j}])\,t^{m-1-j}
\right]_{<m},
\qquad
u_b(t)=\sum_{k\ge0}\frac{\zeta^{(m+k)}(b)}{(m+k)!}t^k.
\tag{DPL5.5}
\] All other primary jets vanish. Indeed \(D\) identifies this dual
generalized eigenspace with the source generalized eigenspace at \(b\),
by (DPL0.4) and (DPL5.2). GZR4 gives \[
\lambda([E_{a,j}])
=(-1)^j[t^{m-1-j}]\bigl(j_bx(t)/u_b(t)\bigr).
\] Solving these \(m\) coefficient equations proves (DPL5.5). Inserting
this polynomial into the existing RZ section at \(b\) produces an actual
element of \(Q\), with its original Mellin representative and all Gamma
and endpoint factors from (DPL1.1), not merely a formal jet.

For \(a\notin\mathscr Z\), the dual generalized eigenspace is zero by
(DPL5.2) and the original resolvent. Repeated factors and all nilpotents
are included in these statements.

\subsection{DPL6. The exact map to the cokernel in the actual ASD14
cone}\label{dpl6.-the-exact-map-to-the-cokernel-in-the-actual-asd14-cone}

Keep every degree and sign of ASD14: \[
D_{\mathrm{full}}=[P\xrightarrow d A_{\mathrm{raw}}],
\qquad
T=\chi\otimes D_{\mathrm{full}}^\vee[-2],
\] \[
K_\zeta=
\left[
P\xrightarrow{-d}A_{\mathrm{raw}}
\xrightarrow{\Psi_\zeta^1}\chi\otimes A_{\mathrm{raw}}'
\xrightarrow{d'}\chi\otimes P'
\right],
\quad\deg=(-1,0,1,2),
\] \[
\Psi_\zeta^1=\pi'D\pi,\qquad
H^1(T)=Y,\qquad H^1(K_\zeta)=\widetilde C.
\tag{DPL6.1}
\] The notation \(A_{\mathrm{raw}}\) here names ASD's original overlap
test space; it is not the transpose operator \(A=L^t\) of DPL0.

The relevant part of the actual comparison triangle's long exact
sequence is \[
0\longrightarrow Q\xrightarrow D Y
\xrightarrow{\eta}\widetilde C
\longrightarrow H^2(D_{\mathrm{full}})=0.
\tag{DPL6.2}
\] Thus the map from the dual cohomology to the cone cohomology is
exactly the quotient map \(\eta\). The following arrow, the connecting
map to \(H^2(D_{\mathrm{full}})\), is zero because its actual codomain
is zero. We do not rename an unspecified arrow a connecting map.

DPL5 proves that \(\eta\) vanishes on every generalized eigenspace of
the full twisted generator \(M=1-L^t\), and more generally on every
polynomial-annihilated submodule of \(Y\). Each such dual class lifts
uniquely to a source cohomology class in \(Q\).

At the actual cochain level, a degree-one cone cycle has the form
\(\pi'\lambda\), because \(\ker d'=\pi'Q'\). If its dual cohomology
class \(\lambda\) satisfies the polynomial condition, write uniquely
\(\lambda=Dx\) by DPL5. Choose any \(b\in A_{\mathrm{raw}}\) with
\(\pi b=x\). Then \[
\pi'\lambda=\pi'D\pi b=\Psi_\zeta^1 b=d_{K_\zeta}^0 b.
\tag{DPL6.3}
\] This is the explicit null-boundary formula, with the positive
middle-cone differential retained. The cohomology lift \(x\) is unique.
The test-space representative \(b\) is determined modulo the actual
summation image \(J=\ker\pi=d(P)\); no unique test-space section has
been asserted.

On \(H^1(K_\zeta)\) itself there are no nonzero generalized
eigenvectors, by DPL3. This conclusion does not remove the retained
groups \[
H^{-1}(K_\zeta)=H,\qquad
H^2(K_\zeta)=\chi\otimes H',
\] or their full faithful extra closed copies, endpoint lines, and
original actions. It does not prove that \(H^1(K_\zeta)=0\) as an
algebraic vector space.

\subsection{DPL7. The actual companion mirror and supported
labels}\label{dpl7.-the-actual-companion-mirror-and-supported-labels}

The mirror of ASD14 compares the original pairing with its companion
denominator \(\zeta(1-s)\). It does not assume the original pairing is
self-mirror. The exact identity is \[
D^\zeta R=-R'D^{\zeta^\vee}.
\tag{DPL7.1}
\] On the degree-one terms in the comparison, the source mirror is
\(-R\) and the target mirror in the actual cone map is \(+R'\).
Therefore a companion lift \(D^{\zeta^\vee}x=\lambda\) transforms into
\[
D^\zeta(-Rx)=R'\lambda.
\tag{DPL7.2}
\] This is the correct signed lift on the original comparison. The
mirror conjugates the generator to \(1\) minus the generator, as follows
either by differentiating ASD14.13 or directly from \(RL=(1-L)R\). It
consequently transports the polynomial \(q(X)\) to \(q(1-X)\). DPL3
applies to both, so polynomial lifting is preserved by this exact
companion mirror, including its signs.

For the existing lattice carrier \[
G_{\mathcal L}(V)=
\{(0,\lambda):\lambda\in\mathcal L\}
\cup\{(v,1_{\mathcal L}):v\in V\},
\] every receiving linear map above lifts by \[
(v,\lambda)\longmapsto(f(v),\lambda).
\tag{DPL7.3}
\] This is well-defined because a nonzero input amplitude has top
support, and a vanishing output retains its original label. Composition
and identities are preserved by substitution. On each polynomial-kernel
subspace the inverse in (DPL5.2) has the same lift, so the unique lift
retains every support label, including supported zero labels. No
ordinary tensor operation identifying different supported zero labels is
substituted for this map.

\subsection{DPL8. Exact consequence and limit of the
calculation}\label{dpl8.-exact-consequence-and-limit-of-the-calculation}

The actual algebraic dual cokernel has no nonzero polynomial-annihilated
vector and no nonzero polynomial-annihilated quotient. The second
statement follows from the reverse Hom vanishing in DPL4 for the
quotient map. Its polynomial operators are continuously invertible in
the stated quotient topologies. Every polynomial-annihilated module has
the two Hom vanishings and the explicitly proved unique extension
splitting in DPL4.

The actual ASD14 map from dual cohomology to its cone cokernel vanishes
on every generalized eigenspace, with a unique full-multiplicity source
lift and the explicit cochain boundary (DPL6.3). This strengthens
density of \(DQ\) to a polynomial lifting theorem. It does not replace
the full dual by its algebraic generalized eigenspaces, assert
convergence of an infinite primary expansion, or conclude that the
entire cokernel vanishes. No Deligne purity bound or RH conclusion has
been inferred from this lifting result.

\subsection{DPL9. Direct cyclic derived-Hom calculation on every degree
of the actual
cone}\label{dpl9.-direct-cyclic-derived-hom-calculation-on-every-degree-of-the-actual-cone}

This final calculation is in the algebraic receiving \(R\)-module
category. It does not assume an analytic derived category or a spectral
sequence. Let \(V=R/(q)\), \(q\ne0\). Its explicit free resolution is \[
P_V^{-1}=R\xrightarrow{\,q\,}P_V^0=R
\longrightarrow V\longrightarrow0.
\tag{DPL9.1}
\] Multiplication by \(q\) is injective because \(R=\mathbb C[X]\) has
no zero divisors, so this is an exact two-term projective resolution.

On the actual cone (DPL6.1) retain the generator \[
G_K^{-1}=L_P,\qquad G_K^0=L_{A_{\mathrm{raw}}},
\qquad G_K^1=1-L_{A_{\mathrm{raw}}}^t,\qquad
G_K^2=1-L_P^t.
\tag{DPL9.2}
\] Here \(L_P,L_{A_{\mathrm{raw}}}\) are the derivatives of the original
real actions on those original spaces, as in ASD7 and ASD14. The
identities \(dL_P=L_{A_{\mathrm{raw}}}d\),
\(\Psi_\zeta^1L_{A_{\mathrm{raw}}}=(1-L_{A_{\mathrm{raw}}}^t)\Psi_\zeta^1\),
and their transpose prove that \(G_K\) commutes with every differential,
including the first minus sign. Thus \(q(G_K)\) is an actual cochain
operator.

The Hom complex of (DPL9.1) into \(K_\zeta\), which computes
\(\operatorname{RHom}_R(V,K_\zeta)\), has terms and differential \[
\mathscr H^n=K_\zeta^n\oplus K_\zeta^{n-1},\qquad
\delta^n(u,v)=\bigl(d_Ku,\ d_Kv-(-1)^nq(G_K)u\bigr).
\tag{DPL9.3}
\] These are exactly the Hom differential \(d_Kf-(-1)^nf\,d_{P_V}\).
Substitution verifies \(\delta^{n+1}\delta^n=0\): the two mixed
polynomial terms have opposite signs.

There is a direct exact cohomology sequence \[
0\longrightarrow
\frac{H^{n-1}(K_\zeta)}{q(G_K)H^{n-1}(K_\zeta)}
\longrightarrow H^n(\mathscr H)
\longrightarrow
\ker\bigl(q(G_K):H^n(K_\zeta)\to H^n(K_\zeta)\bigr)
\longrightarrow0.
\tag{DPL9.4}
\] Here is a proof with all maps. A cocycle \((u,v)\) obeys \(d_Ku=0\)
and \(d_Kv=(-1)^nq(G_K)u\), so its image is the class \([u]\) in the
displayed kernel. Every such class lifts: choose a cycle \(u\)
representing it and \(v\) with the required differential. A class in the
kernel of this map has \(u=d_Ka\). Subtracting the boundary \[
\delta^{n-1}(a,0)=(d_Ka,(-1)^nq(G_K)a)
\] leaves \((0,v-(-1)^nq(G_K)a)\), whose second component is a cycle.
Two such second components differ by a boundary and a polynomial
multiple of a cycle, exactly the left quotient in (DPL9.4). Its
injection sends a cycle class \([v]\) to \([(0,v)]\). This proves all
parts of (DPL9.4) directly.

Use the actual ASD14 cohomology \[
H^{-1}(K_\zeta)=H,\quad H^0(K_\zeta)=0,\quad
H^1(K_\zeta)=\widetilde C,\quad H^2(K_\zeta)=\chi\otimes H',
\] and DPL3's bijectivity of \(q(M)\) on \(\widetilde C\). Equation
(DPL9.4) then yields every degree: \[
\boxed{
\begin{aligned}
H^{-1}\operatorname{RHom}_R(V,K_\zeta)
&=\ker(q(L_H):H\to H),\\
H^{0}\operatorname{RHom}_R(V,K_\zeta)
&=H/q(L_H)H,\\
H^{1}\operatorname{RHom}_R(V,K_\zeta)
&=0,\\
H^{2}\operatorname{RHom}_R(V,K_\zeta)
&=\ker(q(1-L_H^t):\chi\otimes H'\to\chi\otimes H'),\\
H^{3}\operatorname{RHom}_R(V,K_\zeta)
&=(\chi\otimes H')/q(1-L_H^t)(\chi\otimes H'),\\
H^{n}\operatorname{RHom}_R(V,K_\zeta)&=0
\quad(n\notin\{-1,0,1,2,3\}).
\end{aligned}}
\tag{DPL9.5}
\] The induced \(L_H\) is the original action on \(H=\ker d\). In
particular \[
H^2\operatorname{RHom}_R(V,K_\zeta)
=\operatorname{Hom}_R(V,\chi\otimes H').
\tag{DPL9.6}
\] The identity follows by evaluating an \(R\)-linear map on the class
of \(1\) in \(R/(q)\). Likewise the degree-zero and degree-three
quotients are the ordinary cyclic \(\operatorname{Ext}^1_R\) terms
computed by (DPL9.1). No purity assertion is involved.

The potentially surviving terms in (DPL9.5) retain \(H,H'\), their four
original endpoint lines, and both faithful extra closed coefficient
copies. The result proves the specified degree-one vanishing and
calculates the exact other terms; it does not assert that the entire
comparison is a quasi-isomorphism or that any of the retained terms
obstructs the user's full source geometry.

\subsection{Exact original character lifting and its full supported
return}\label{exact-original-character-lifting-and-its-full-supported-return-occurrence-2}

MCL0--MCL11 now computes the original transpose A′→S′ on every actual
zero and all its Mellin jets. A target jet of order j needs
generalized-character length exactly m+j+1 at a zero of multiplicity m.
The full Euler-distribution classification proves minimality against the
whole original continuous dual; two actual primes remove the
single-prime character aliases. Every prime logarithm and the full
Gamma/Fourier return remain. This is an actual continuous lift, with a
nonzero class only when its original annihilator must be preserved.

RPC0--RPC8 constructs the original restriction's actual pushout by the
residue image, proves its precise cochain location and a single coherent
real/prime-equivariant section on all locally finite classes.
ORE0--ORE10 calculates the unmodified connecting class, including its
full Gamma germ and rank min(m,r), then proves its place in every degree
of the full supported derived comparison. The degree-two map is
(eta,e)↦(−delta\_q(eta)/2,e). Thus its image in C\_zeta/qC\_zeta is
zero, while the full supported map retains its Y/qY component. Both
assertions follow from the constructed maps; no cone degree or extra
coefficient is discarded.

These results refine the meaning of the prior finite-lift statement
without withdrawing its proved scope. Fixed-annihilator splitting is
stronger than the surjectivity in Deligne's invariant-cycle theorem. The
original longer-chain map is already surjective. The present numerical
character truncation retains its same-character extension whenever it
retains the target; it is not Deligne's independently proved geometric
weight gap. Complete proofs, independent checks and inspected
reproducible diagrams have been inserted into the cumulative TeX. The
actual tau weight-separation goal remains active.
